\subsection*{Sprint 7: Persistencia en la nube}
\addcontentsline{toc}{subsection}{Sprint 7: Persistencia en la nube}

El séptimo sprint completa la \textbf{persistencia de datos} de
\textbf{Muncher} llevando a la nube los mecanismos establecidos en el sprint
anterior. Al inicio del sprint, la API dispone en local de una base de datos
relacional, obtiene sus datos de acceso de variables de entorno y cuenta con
un gestor de migraciones que las pruebas de extremo a extremo ya ejercitan en
cada ejecución. Los entornos de desarrollo y de producción, en cambio, siguen
sin base de datos, de modo que ninguna funcionalidad que dependa de ella podría
desplegarse todavía.

El primer objetivo es desplegar una \textbf{base de datos relacional} en cada
entorno de la nube, declarada en la misma plantilla que el resto de la
infraestructura para que se cree, se actualice y se destruya junto con el
entorno al que pertenece. La contraseña de la base de datos se generará
automáticamente y se almacenará en un \textbf{secreto}, de forma que no
aparezca ni en el código ni en la plantilla, tal como exige
\reqref{RNF-SE-03}. Los datos de acceso se inyectarán en las variables de
entorno de la API, que es exactamente el mecanismo por el que la API se conecta
ya en local: la base de datos de la nube no exige, por tanto, ningún cambio en
el código de la aplicación.

El segundo objetivo es que las migraciones se ejecuten como parte del
\textbf{despliegue continuo}. Si el esquema de un entorno se actualizara de
forma manual, nada garantizaría que corresponda con la versión de la API que se
ejecuta en él, y una migración olvidada se manifestaría como un fallo en
producción. Al incorporarlas al ciclo de despliegue, cada versión del código
llega al entorno acompañada del esquema que necesita, y un fallo al aplicarlas
detiene el despliegue antes de que la nueva versión atienda peticiones.

Por último, la monitorización establecida en el quinto sprint se extenderá a
la base de datos. Igual que entonces, el propósito es que cada componente nuevo
nazca ya instrumentado: la base de datos es el primer componente del sistema
que conserva estado, y los problemas que pueden originarse en ella —consultas
lentas, conexiones agotadas o falta de espacio— no se manifiestan
necesariamente como errores de la API, sino como respuestas que se degradan.
Observarla desde el principio permite distinguir el tiempo que la API dedica a
la base de datos del que dedica a su propio procesamiento.

\subsubsection*{Objetivos iniciales}

\begin{itemize}
	\item Seleccionar el servicio de base de datos gestionado de la nube.
	\item Declarar la base de datos en la plantilla de infraestructura, generando
	      su contraseña en un secreto e inyectando los datos de acceso en las
	      variables de entorno de la API.
	\item Ejecutar las migraciones como parte del ciclo de despliegue continuo.
	\item Extender la monitorización para que observe también la base de datos.
	\item Desplegar la base de datos y la API conectada a ella en el entorno de
	      desarrollo.
	\item Desplegar la base de datos y la API conectada a ella en el entorno de
	      producción.
\end{itemize}

\subsubsection*{Historias de usuario completadas}

\begin{center}
	\begin{tabularx}{\textwidth}{|l|X|c|}
		\hline
		\textbf{Identificador} & \textbf{Nombre}                       & \textbf{Puntos de historia} \\
		\hline
		\usref{US-DE-06}       & Base de datos en la nube              & 3                           \\
		\hline
		\usref{US-DE-07}       & Migraciones en el despliegue continuo & 2                           \\
		\hline
		\usref{US-DE-08}       & Monitorización de la base de datos    & 3                           \\
		\hline
	\end{tabularx}
\end{center}
